\section{Discussion}

The main empirical risk is overinterpreting standard-split residual gains. A residual model can improve forecasting by reading local event evidence, but it can also learn that certain positions and velocities imply a likely future acceleration under the dataset collection protocol. The linear residual diagnostics and decorrelation experiments are intended to separate these explanations. If a linear model over filtered state improves substantially, then some part of the residual signal is available without images. If image-conditioned gains survive decorrelation and exceed the filter-state-only residual, then the evidence for visual conditioning is stronger.

The decorrelated subset should not be interpreted as a true out-of-distribution benchmark. It removes a specified class of linear motion-prior shortcuts after sample capping, and the greedy candidate search is approximate. Nonlinear shortcuts, remaining split-specific behavior, and visual context cues can persist. The optional mean-acceleration penalty addresses a global acceleration bias, such as a downward tendency, but choosing its weight is application dependent: for a horizontal camera it may be a dataset artifact, while for some deployment settings gravity-related image-plane acceleration may be a real prior.

The preliminary single-input representation results also require caution. Event frames perform well on the raw split, while CSTR2 appears stronger in early decorrelated runs. One possible explanation is that CSTR2 preserves local motion direction through timestamp gradients, whereas event frames are closer to binary occupancy. However, this has not yet been isolated by controlled ablations over temporal bin count, crop size, or backbone capacity.

The strongest claim is therefore not that one event representation is universally best. Instead, the results should be interpreted as evidence that (i) full-box optimized Kalman filtering is a strong and necessary baseline on FRED, (ii) single-input event-conditioned residuals can improve box forecasting under the standard protocol, and (iii) correlation analysis and decorrelation are necessary diagnostics before attributing those gains to visual evidence or OOD-ready dynamics.
